\documentclass[11pt]{article}
\usepackage[margin=1in]{geometry}
\usepackage{amsmath,amssymb,amsthm}
\usepackage{graphicx}
\usepackage[hidelinks]{hyperref}

\newtheorem*{theorem}{Theorem}

\title{A Renorming Answer for the Schreier Space and the Uniform \texorpdfstring{$\lambda$}{lambda}-Property}
\author{Run fa\_banach\_001}
\date{2026-06-20}

\begin{document}
\maketitle

\section*{Status}

This packet records a literature-implied full answer to Question (II) in Section 5 of Beanland--Chu, \emph{The Schreier Space Does Not Have the Uniform $\lambda$-property}, arXiv:1904.10384. The question asks whether the Schreier space $X$ can be renormed to have the uniform $\lambda$-property. The answer is yes.

\section*{Source Question}

The source question appears on page 7, Section 5, item (II).

\begin{center}
\includegraphics[width=0.98\linewidth]{figures/open_problem_crop.png}
\end{center}

The relevant item is:
\[
\text{(II) Can } X \text{ be renormed to have the uniform } \lambda\text{-property?}
\]

\section*{Identification}

The Schreier space is separable because it is the completion of $c_{00}$ under a norm for which the standard finitely supported rational vectors are dense. Every separable Banach space admits an equivalent strictly convex norm. For completeness, the standard short construction is recalled in the proof below. Once the norm is strictly convex, the unit sphere is contained in the set of extreme points of the unit ball. This immediately gives the uniform $\lambda$-property with constant $1/2$.

\begin{theorem}
Let $X$ be any separable real Banach space. Then $X$ admits an equivalent norm with the uniform $\lambda$-property. In particular, the Schreier space admits such a renorming.
\end{theorem}

\begin{proof}
Let $\|\cdot\|$ be the original norm on a separable Banach space $X$. Choose a countable norm-dense set $(u_n)$ in the unit sphere. By Hahn--Banach, choose $f_n\in X^*$ with $\|f_n\|=1$ and $f_n(u_n)=1$. The family $(f_n)$ is total: if $x\ne 0$, choose $u_n$ with $\|u_n-x/\|x\|\|<1/2$, and then $f_n(x)\ne 0$.

Define
\[
  \|x\|'=\left(\|x\|^2+\sum_{n=1}^{\infty}2^{-n}|f_n(x)|^2\right)^{1/2}.
\]
This norm is equivalent to $\|\cdot\|$, since
\[
  \|x\|\leq \|x\|'\leq \sqrt{2}\,\|x\|.
\]
It is strictly convex. Indeed, equality in the triangle inequality for $\|\cdot\|'$ forces equality in the Hilbert-space norm of the scalar coordinate vectors
\[
  x\mapsto \bigl(\|x\|,2^{-1/2}f_1(x),2^{-2/2}f_2(x),\ldots\bigr).
\]
Thus, for nonzero $x,y$, equality implies that all coordinates are positively proportional, so $f_n(x)=c f_n(y)$ for every $n$. Since $(f_n)$ is total, $x=cy$. Hence two distinct points of the $\|\cdot\|'$-unit sphere cannot have a nontrivial midpoint on the unit sphere.

Now work in the renormed space $(X,\|\cdot\|')$. Strict convexity implies that every unit vector is an extreme point of the closed unit ball. If $x$ is in the closed unit ball and $x\ne 0$, put $e=x/\|x\|'$ and $t=\|x\|'$. Then
\[
  x=t e=\frac{1+t}{2}e+\frac{1-t}{2}(-e).
\]
Here $e$ is an extreme point, $-e$ is in the unit ball, and the coefficient $(1+t)/2$ is at least $1/2$. For $x=0$, the same formula with any unit vector $e$ gives
\[
  0=\frac12 e+\frac12(-e).
\]
Therefore the renormed space has the uniform $\lambda$-property with constant $\lambda_0=1/2$.
\end{proof}

\section*{Scope and Limitations}

This answers only the renorming question, item (II), for equivalent renormings. It does not settle the source paper's item (I), which asks whether higher-order Schreier spaces and their duals fail the uniform $\lambda$-property in their canonical norms, nor item (III), which asks about sharpening the $1/8$ bound for $p$-convexifications.

\section*{Novelty Check}

Searches on 2026-06-20 for exact phrases around ``Can $X$ be renormed'', ``Schreier space renormed uniform lambda-property'', and ``strictly convex uniform lambda-property'' did not find a later Schreier-specific paper explicitly answering item (II). The underlying renorming mechanism is classical and standard; the contribution of this packet is the direct identification that it answers the source question.

\begin{thebibliography}{9}
\bibitem{BeanlandChu}
K. Beanland and H. V. Chu,
\emph{The Schreier Space Does Not Have the Uniform $\lambda$-property},
arXiv:1904.10384, 2019.

\bibitem{AronLohman}
R. M. Aron and R. H. Lohman,
\emph{A geometric function determined by extreme points of the unit ball of a normed space},
Pacific J. Math. 127 (1987), no. 2, 209--231.
\end{thebibliography}

\end{document}
